% >>> typage-template springer-lncs@2026.10.1 variant=runningheads
\documentclass[runningheads]{llncs}
\bibliographystyle{splncs04}
% <<< typage-template
\usepackage{amsmath}
\usepackage{graphicx}
\usepackage{booktabs}

\begin{document}

\title{Contribution Title\thanks{Supported by Funding Organization.}}

\author{First Author\inst{1} \and Second Author\inst{2}}

\authorrunning{F. Author and S. Author}

\institute{Department, University, City, Country\\
\email{first.author@example.com}
\and
Laboratory, Institute, City, Country\\
\email{second.author@example.com}}

\maketitle

\begin{abstract}
The abstract should briefly summarize the contents of the paper in
150--250 words.

\keywords{First keyword \and Second keyword \and Third keyword.}
\end{abstract}

\section{Introduction}
Introduce the problem and list the contributions. Cite prior work by
number~\cite{lamport1994latex,knuth1984texbook}.

\section{Method}
Describe the method. Number the equations you refer to:
\begin{equation}
  \label{eq:objective}
  \min_{\theta} \; \frac{1}{n} \sum_{i=1}^{n} \ell\bigl(f_\theta(x_i), y_i\bigr) .
\end{equation}

\begin{theorem}
State a result in the theorem environment provided by the class.
\end{theorem}

\begin{proof}
Prove it.
\end{proof}

\section{Experiments}
Table~\ref{tab:results} summarizes the results of Eq.~\eqref{eq:objective}.

\begin{table}
  \caption{Results on the example benchmark.}
  \label{tab:results}
  \centering
  \begin{tabular}{lcc}
    \toprule
    Method   & Accuracy (\%) & Time (s) \\
    \midrule
    Baseline & 81.2          & 12.4     \\
    Ours     & 84.7          & 10.9     \\
    \bottomrule
  \end{tabular}
\end{table}

\section{Conclusion}
Conclude and outline future work.

\begin{credits}
\subsubsection{\ackname} Acknowledge funding and help here.

\subsubsection{\discintname}
The authors have no competing interests to declare that are relevant to the
content of this article.
\end{credits}

\bibliography{refs}

\end{document}
